\chapter{Appendices}
\label{ch:appendix}
\renewcommand{\thesection}{\Alph{section}}
\setcounter{section}{0}

\section{Hyperparameter Inventory}
\label{app:hparams}
Table~\ref{tab:hparams} lists the main hyperparameters, gathered in one place so a reader can reproduce
any cell. Classical and BiLSTM settings are the defaults used throughout; the Transformer and LLM
settings follow standard fine-tuning practice with validation-selected learning rates
(\S\ref{sec:training}). The text axes (the preprocessing mode and input format, \S\ref{sec:axes}) and the
experiment registry (\S\ref{sec:grid}) define which combinations of these settings each experiment
sweeps; the exact LLM prompt template is given in \S\ref{sec:pillars}.

\begin{table}[ht]
\centering
\caption{Key hyperparameters of each model.}
\label{tab:hparams}
\small
\begin{tabular}{@{}ll>{\raggedright\arraybackslash}p{7.5cm}@{}}
\toprule
Pillar & Setting & Value \\
\midrule
Classical   & TF-IDF analyzer / $n$-grams & \texttt{char\_wb}, $(2,4)$ \\
            & max features / SVR $C$ / kernel & 15{,}000 / 1.0 / RBF \\
            & extra baselines & \texttt{tfidf\_cos}, \texttt{fasttext\_cos} (100-dim) \\
RNN         & char vocab / embed / hidden / layers & 5{,}000 / 128 / 128 / 2 (bidirectional) \\
            & dropout / batch / max epochs & 0.3 / 64 / 20 (patience 4) \\
            & learning rate / optimiser & $1\mathrm{e}{-}3$ / Adam \\
Transformer & backbones & mBERT, XLM-R, GTE-multilingual \\
            & config / frozen layers / max length / batch & dual + cross / 6 / 256 / 16 \\
            & learning rate / weight decay / dropout & $2\mathrm{e}{-}5$ / 0.01 / 0.2 \\
LLM         & base / adaptation & Qwen 3.5 4B / QLoRA (4-bit NF4, double-quant) \\
            & LoRA rank / alpha / dropout & 16 / 16 / 0 \\
            & target modules & q, k, v, o, gate, up, down \\
            & epochs / batch (effective) / warm-up & 10 / 4 (16 via accumulation) / 5\,\% \\
            & max sequence length / loss & 1{,}024 / completion-only \\
            & learning rate (val-selected) & $\in\{1,2,3\}\times10^{-4}$ \\
\midrule
Common      & preprocessing modes / input formats & \{raw, clean, segment\} / \{ra, qar\} \\
            & max-score feature normaliser & divide by 20 (largest observed max) \\
            & split / seed & 70/15/15 / 42 (multi-seed: 42, 1337, 2024, 7, 99) \\
            & target / loss & $\hat{y}\in[0,1]$ / mean-squared error (regression heads) \\
\bottomrule
\end{tabular}
\end{table}

\section{Compute Budget}
\label{app:compute}
The classical models train in seconds on a CPU. The full neural, Transformer, and LLM grid (hundreds
of cells across two dataset variants) requires roughly 40 GPU-hours in total. All metrics are
computed from saved per-answer predictions and need no GPU.

\section{Full Experiment Grid}
\label{app:grid}
The complete per-cell leaderboards, every combination of model, preprocessing mode (\texttt{raw} /
\texttt{clean} / \texttt{segment}), input format (\texttt{ra} / \texttt{qar}), and dataset variant
(\texttt{full} / \texttt{no10c}), with train and test metrics, are released with the code under
\texttt{results/leaderboards/}. Chapter~\ref{ch:results} reports the validation-selected champion per
family and the one-factor-at-a-time effect of each axis (\S\ref{sec:results-data}); the exhaustive grid
is kept here to keep the chapter readable.

\section{Frontier-Model Breakdown}
\label{app:frontier}
Table~\ref{tab:frontier} in Chapter~\ref{ch:results} reports each frontier model at its best of the
two zero-shot settings. Table~\ref{tab:frontier-full} gives the full per-setting breakdown (bare
integer vs.\ reasoning-enabled) on the \texttt{no10c} test set, with the per-model unparsed rate;
unparseable responses are excluded from the metrics rather than imputed. Each model version and the
evaluation date are pinned, and the raw API responses are cached so the snapshot reproduces.

\begin{table}[ht]
\centering
\caption{Frontier models, full zero-shot breakdown (bare vs.\ reasoning). ``[API-pending]'' awaits the
run; \texttt{exp14} writes \texttt{results\_frontier/no10c/frontier\_metrics.csv}.}
\label{tab:frontier-full}
\small
\begin{tabular}{@{}llccc@{}}
\toprule
Model & Setting & QWK & Exact & Unparsed \\
\midrule
GPT      & bare      & [API-pending] & [API-pending] & [API-pending] \\
GPT      & reasoning & [API-pending] & [API-pending] & [API-pending] \\
Claude   & bare      & [API-pending] & [API-pending] & [API-pending] \\
Claude   & reasoning & [API-pending] & [API-pending] & [API-pending] \\
Gemini   & bare      & [API-pending] & [API-pending] & [API-pending] \\
Gemini   & reasoning & [API-pending] & [API-pending] & [API-pending] \\
DeepSeek & bare      & [API-pending] & [API-pending] & [API-pending] \\
DeepSeek & reasoning & [API-pending] & [API-pending] & [API-pending] \\
\bottomrule
\end{tabular}
\end{table}

\section{Additional Explanation Heatmaps}
\label{app:heatmaps}
Chapter~\ref{ch:results} shows the classical SHAP heatmap (Figure~\ref{fig:heatmap}). For
completeness, Figures~\ref{fig:heatmap-bilstm}, \ref{fig:heatmap-encoder}, and \ref{fig:heatmap-llm}
show SHAP heatmaps for the remaining three pillars on real Khmer answers, so all four
families are illustrated. Khmer complex script is rendered with browser text shaping, as standard
plotting back-ends do not correctly stack subscript consonants or reorder vowels, and the readable
tokens are used.

\begin{figure}[ht]
\centering
\includegraphics[width=0.92\textwidth]{heatmap_bilstm.png}
\caption{Sample of a Khmer word-attribution heatmap (BiLSTM model, SHAP).}
\label{fig:heatmap-bilstm}
\end{figure}

\begin{figure}[ht]
\centering
\includegraphics[width=0.92\textwidth]{heatmap_encoder.png}
\caption{Sample of a Khmer word-attribution heatmap (Transformer model, SHAP).}
\label{fig:heatmap-encoder}
\end{figure}

\begin{figure}[ht]
\centering
\includegraphics[width=0.92\textwidth]{heatmap_llm.png}
\caption{Sample of a Khmer word-attribution heatmap (LLM model, SHAP).}
\label{fig:heatmap-llm}
\end{figure}

\section{Live Prototype Screenshots}
\label{app:prototype}
\emph{[Insert here one or two screenshots of the live prototype: the input form (Question + Reference +
Answer + model selector), and a returned score with its Khmer word-attribution heatmap and written
feedback.]}

\section{Ethics and Data Governance}
\label{app:ethics}
The corpus contains real, anonymised student answers from minors, collected from two Cambodian high
schools and used solely for research into a teacher-assist grading tool. Access to the examination
scripts was granted by approaching each school and obtaining permission \emph{directly from the school
Director}; no formal institutional review board (IRB) process applies at this level in this setting.
Only pseudonymous identifiers are stored; no student names appear in the data. The system is intended
to \emph{support}, not replace, a teacher's judgement (human-in-the-loop). Because a single teacher
graded all answers, the models reflect that grader's decisions; results are reported as agreement with
this grader, with reporting across two dataset variants as partial mitigation. The grader
inconsistency documented in \S\ref{sec:noisy} (class 10C Biology) is disclosed rather than hidden.
\emph{[If the institution later requires formal consent/IRB documentation, insert the reference
here.]}
